\documentclass{article}
\usepackage[utf8]{inputenc}
\usepackage[spanish]{babel}
\usepackage{amsmath, amssymb}

\begin{document}

\section*{Resolución de ecuaciones irracionales}

\subsection*{a) $x - \sqrt{2x - 3} = 1$}
\begin{align*}
x - 1 &= \sqrt{2x - 3} && \text{(Aislamos la raíz)} \\
(x - 1)^2 &= \left(\sqrt{2x - 3}\right)^2 && \text{(Elevamos ambos miembros al cuadrado)} \\
x^2 - 2x + 1 &= 2x - 3 \\
x^2 - 4x + 4 &= 0 \\
(x - 2)^2 &= 0 \implies x = 2
\end{align*}
\textbf{Comprobación:}
\[
2 - \sqrt{2(2) - 3} = 2 - \sqrt{1} = 2 - 1 = 1 \quad (\text{Válida})
\]
\textbf{Solución:} $x = 2$

\hrule

\subsection*{b) $\sqrt{x + 4} - \sqrt{6 - x} = -2$}
\begin{align*}
\sqrt{x + 4} &= \sqrt{6 - x} - 2 && \text{(Aislamos una de las raíces)} \\
\left(\sqrt{x + 4}\right)^2 &= \left(\sqrt{6 - x} - 2\right)^2 && \text{(Elevamos al cuadrado)} \\
x + 4 &= (6 - x) - 4\sqrt{6 - x} + 4 \\
x + 4 &= 10 - x - 4\sqrt{6 - x} \\
4\sqrt{6 - x} &= 6 - 2x && \text{(Aislamos la raíz restante)} \\
2\sqrt{6 - x} &= 3 - x && \text{(Dividimos entre 2 para simplificar)} \\
\left(2\sqrt{6 - x}\right)^2 &= (3 - x)^2 && \text{(Elevamos al cuadrado de nuevo)} \\
4(6 - x) &= 9 - 6x + x^2 \\
24 - 4x &= 9 - 6x + x^2 \\
x^2 - 2x - 15 &= 0
\end{align*}
Aplicamos la fórmula cuadrática:
\[
x = \frac{2 \pm \sqrt{(-2)^2 - 4(1)(-15)}}{2(1)} = \frac{2 \pm \sqrt{64}}{2} = \frac{2 \pm 8}{2} \implies 
\begin{cases} 
x_1 = 5 \\ 
x_2 = -3 
\end{cases}
\]
\textbf{Comprobación:}
\begin{itemize}
    \item Para $x = 5$: $\sqrt{5 + 4} - \sqrt{6 - 5} = \sqrt{9} - \sqrt{1} = 3 - 1 = 2 \neq -2$ (No válida)
    \item Para $x = -3$: $\sqrt{-3 + 4} - \sqrt{6 - (-3)} = \sqrt{1} - \sqrt{9} = 1 - 3 = -2$ (Válida)
\end{itemize}
\textbf{Solución:} $x = -3$

\hrule

\subsection*{c) $\sqrt{x^2 + 2x + 9} - 7 = 2x$}
\begin{align*}
\sqrt{x^2 + 2x + 9} &= 2x + 7 && \text{(Aislamos la raíz)} \\
\left(\sqrt{x^2 + 2x + 9}\right)^2 &= (2x + 7)^2 && \text{(Elevamos al cuadrado)} \\
x^2 + 2x + 9 &= 4x^2 + 28x + 49 \\
3x^2 + 26x + 40 &= 0
\end{align*}
Aplicamos la fórmula cuadrática:
\[
x = \frac{-26 \pm \sqrt{26^2 - 4(3)(40)}}{2(3)} = \frac{-26 \pm \sqrt{676 - 480}}{6} = \frac{-26 \pm \sqrt{196}}{6} = \frac{-26 \pm 14}{6}
\]
De donde obtenemos:
\[
x_1 = \frac{-12}{6} = -2, \qquad x_2 = \frac{-40}{6} = -\frac{20}{3}
\]
\textbf{Comprobación:}
\begin{itemize}
    \item Para $x = -2$: $\sqrt{(-2)^2 + 2(-2) + 9} - 7 = \sqrt{9} - 7 = 3 - 7 = -4 = 2(-2)$ (Válida)
    \item Para $x = -\frac{20}{3}$: $\sqrt{\left(-\frac{20}{3}\right)^2 + 2\left(-\frac{20}{3}\right) + 9} - 7 = \sqrt{\frac{361}{9}} - 7 = \frac{19}{3} - 7 = -\frac{2}{3} \neq -\frac{40}{3}$ (No válida)
\end{itemize}
\textbf{Solución:} $x = -2$

\hrule

\subsection*{d) $\sqrt{20 - x} = x - 8$}
\begin{align*}
\left(\sqrt{20 - x}\right)^2 &= (x - 8)^2 && \text{(Elevamos ambos miembros al cuadrado)} \\
20 - x &= x^2 - 16x + 64 \\
x^2 - 15x + 44 &= 0 \\
(x - 11)(x - 4) &= 0 \implies 
\begin{cases} 
x_1 = 11 \\ 
x_2 = 4 
\end{cases}
\end{align*}
\textbf{Comprobación:}
\begin{itemize}
    \item Para $x = 11$: $\sqrt{20 - 11} = \sqrt{9} = 3$, y $11 - 8 = 3 \implies 3 = 3$ (Válida)
    \item Para $x = 4$: $\sqrt{20 - 4} = \sqrt{16} = 4$, y $4 - 8 = -4 \implies 4 \neq -4$ (No válida)
\end{itemize}
\textbf{Solución:} $x = 11$

\end{document}